\renewcommand{\tema}{Examen Diagnóstico - Óptica}

\twocolumn
\section*{PREGUNTAS}

\begin{preguntabox}
\subsection*{Pregunta 1}
La luz visible es una onda electromagnética. En cuanto a su tipo de vibración, la luz es:

\begin{enumerate}
    \item[A)] Longitudinal, porque se propaga en línea recta.
    \item[B)] Longitudinal, porque no requiere medio material para propagarse.
    \item[C)] Transversal, porque los campos eléctrico y magnético oscilan perpendicularmente a la dirección de propagación.
    \item[D)] Transversal, porque viaja más lento en medios densos.
\end{enumerate}
\end{preguntabox}

\begin{preguntabox}
\subsection*{Pregunta 2}
¿Cuál de los siguientes modelos de la luz explica correctamente los fenómenos de difracción e interferencia?

\begin{enumerate}
    \item[A)] Modelo corpuscular de Newton, porque las partículas se desvían alrededor de obstáculos.
    \item[B)] Modelo ondulatorio de Huygens, porque cada punto del frente de onda actúa como fuente secundaria.
    \item[C)] Modelo corpuscular de Newton, porque predice que la luz va más rápido en medios densos.
    \item[D)] Ambos modelos explican igualmente bien la difracción.
\end{enumerate}
\end{preguntabox}

\begin{preguntabox}
\subsection*{Pregunta 3}
Un rayo de luz incide sobre un espejo plano formando un ángulo de $40^\circ$ con la superficie del espejo. ¿Cuál es el ángulo de reflexión?

\begin{enumerate}
    \item[A)] $40^\circ$
    \item[B)] $50^\circ$
    \item[C)] $80^\circ$
    \item[D)] $90^\circ$
\end{enumerate}
\end{preguntabox}
\newpage
\begin{preguntabox}
\subsection*{Pregunta 4}
La velocidad de la luz en cierto vidrio es $2 \times 10^8$ m/s. ¿Cuál es el índice de refracción de ese vidrio?

\begin{enumerate}
    \item[A)] 2.0
    \item[B)] 0.67
    \item[C)] 1.5
    \item[D)] 1.33
\end{enumerate}
\end{preguntabox}

\begin{preguntabox}
\subsection*{Pregunta 5}
Un rayo de luz pasa del aire ($n_1 = 1$) al agua ($n_2 = 1.33$) con un ángulo de incidencia de $30^\circ$. ¿Cuál es el valor de $\sin\theta_2$?

\begin{enumerate}
    \item[A)] 0.665
    \item[B)] 0.376
    \item[C)] 1.33
    \item[D)] 0.500
\end{enumerate}
\end{preguntabox}

\begin{preguntabox}
\subsection*{Pregunta 6}
Un objeto se coloca a 60 cm frente a un espejo cóncavo cuya distancia focal es 20 cm. ¿A qué distancia del espejo se forma la imagen?

\begin{enumerate}
    \item[A)] 15 cm
    \item[B)] 80 cm
    \item[C)] 30 cm
    \item[D)] 40 cm
\end{enumerate}
\end{preguntabox}

\begin{preguntabox}
\subsection*{Pregunta 7}
Para el espejo del problema anterior ($d_o = 60$ cm, $d_i = 30$ cm), ¿cuál es el aumento lateral y cómo es la imagen?

\begin{enumerate}
    \item[A)] $m = +0.5$; imagen derecha y reducida.
    \item[B)] $m = -2$; imagen invertida y amplificada.
    \item[C)] $m = -0.5$; imagen invertida y reducida.
    \item[D)] $m = +2$; imagen derecha y amplificada.
\end{enumerate}
\end{preguntabox}
\newpage
\begin{preguntabox}
\subsection*{Pregunta 8}
Una lente convergente tiene distancia focal de 12 cm. Un objeto se coloca a 36 cm de la lente. ¿A qué distancia se forma la imagen?

\begin{enumerate}
    \item[A)] 24 cm
    \item[B)] 9 cm
    \item[C)] 18 cm
    \item[D)] 48 cm
\end{enumerate}
\end{preguntabox}

\begin{preguntabox}
\subsection*{Pregunta 9}
Una lente convergente de distancia focal 15 cm tiene un objeto colocado a 10 cm (entre el foco y la lente). ¿Dónde se forma la imagen y de qué tipo es?

\begin{enumerate}
    \item[A)] A 30 cm al otro lado de la lente; real e invertida.
    \item[B)] A 6 cm al mismo lado del objeto; virtual y reducida.
    \item[C)] A 30 cm al mismo lado del objeto; virtual y amplificada.
    \item[D)] No se forma imagen porque el objeto está entre $F$ y la lente.
\end{enumerate}
\end{preguntabox}

\begin{preguntabox}
\subsection*{Pregunta 10}
Una lente divergente tiene distancia focal de $-20$ cm. Un objeto se coloca a 60 cm. ¿A qué distancia se forma la imagen?

\begin{enumerate}
    \item[A)] $-15$ cm (virtual, mismo lado que el objeto)
    \item[B)] $+15$ cm (real, lado opuesto al objeto)
    \item[C)] $-30$ cm (virtual, mismo lado que el objeto)
    \item[D)] $+60$ cm (real, lado opuesto al objeto)
\end{enumerate}
\end{preguntabox}

% -------------------------------------------------------
%  RESPUESTAS Y RETROALIMENTACIÓN
% -------------------------------------------------------

\newpage
\onecolumn
\section*{RESPUESTAS Y RETROALIMENTACIÓN}

\begin{respuestabox}
\subsection*{Pregunta 1}
\textcolor{verdecorrecto}{\textbf{Respuesta correcta: C) Transversal, porque los campos eléctrico y magnético oscilan perpendicularmente a la propagación.}}

\textbf{Justificación:}\\
La luz es una onda electromagnética transversal: sus campos $\vec{E}$ y $\vec{B}$ oscilan en planos perpendiculares a la dirección de propagación. Esto la distingue de las ondas longitudinales (como el sonido), donde la perturbación es paralela a la propagación.

\textbf{Respuestas incorrectas:}
\begin{itemize}
    \item \textbf{A):} Propagarse en línea recta no define el tipo de onda; las ondas longitudinales también viajan en línea recta.
    \item \textbf{B):} No requerir medio material es una característica de las ondas electromagnéticas en general, pero no define si son transversales o longitudinales.
    \item \textbf{D):} Viajar más lento en medios densos es consecuencia de la refracción, no del tipo de vibración.
\end{itemize}
\end{respuestabox}

\begin{respuestabox}
\subsection*{Pregunta 2}
\textcolor{verdecorrecto}{\textbf{Respuesta correcta: B) Modelo ondulatorio de Huygens.}}

\textbf{Justificación:}\\
La difracción (doblamiento alrededor de obstáculos) y la interferencia (suma algebraica de amplitudes) son fenómenos exclusivos de las ondas. El modelo corpuscular de Newton no puede explicarlos. El principio de Huygens — cada punto del frente de onda es fuente secundaria — predice naturalmente ambos fenómenos.

\textbf{Respuestas incorrectas:}
\begin{itemize}
    \item \textbf{A):} Las partículas del modelo de Newton no se doblan alrededor de obstáculos; viajan en línea recta.
    \item \textbf{C):} El modelo de Newton predice que la luz va \textit{más rápido} en medios densos — lo cual es incorrecto experimentalmente.
    \item \textbf{D):} Solo el modelo ondulatorio explica difracción e interferencia; el corpuscular no.
\end{itemize}
\end{respuestabox}

\begin{respuestabox}
\subsection*{Pregunta 3}
\textcolor{verdecorrecto}{\textbf{Respuesta correcta: B) $50^\circ$}}

\textbf{Justificación:}\\
Es importante recordar que los ángulos de incidencia y reflexión se miden respecto a la \textbf{normal}, no respecto a la superficie. Si el rayo forma $40^\circ$ con la \textit{superficie}, entonces forma $90^\circ - 40^\circ = 50^\circ$ con la normal.

$$\theta_i = 90^\circ - 40^\circ = 50^\circ \quad \Rightarrow \quad \theta_r = \theta_i = 50^\circ$$

\textbf{Respuestas incorrectas:}
\begin{itemize}
    \item \textbf{A) $40^\circ$:} Error de confundir el ángulo con la superficie con el ángulo con la normal.
    \item \textbf{C) $80^\circ$:} Error de duplicar el ángulo dado: $2 \times 40^\circ$.
    \item \textbf{D) $90^\circ$:} No hay ninguna razón física para que el ángulo de reflexión sea siempre $90^\circ$.
\end{itemize}
\end{respuestabox}
\newpage
\begin{respuestabox}
\subsection*{Pregunta 4}
\textcolor{verdecorrecto}{\textbf{Respuesta correcta: C) 1.5}}

\textbf{Fórmula utilizada:} $n = \dfrac{c}{v}$

\textbf{Sustitución:}
$$n = \frac{3 \times 10^8 \text{ m/s}}{2 \times 10^8 \text{ m/s}} = 1.5$$

\textbf{Respuestas incorrectas:}
\begin{itemize}
    \item \textbf{A) 2.0:} Error de calcular $c \div (c - v)$ o simplemente dividir $3/1.5$ sin justificación.
    \item \textbf{B) 0.67:} Error de invertir el cociente: $v/c = 2/3 \approx 0.67$. El índice siempre es $\geq 1$.
    \item \textbf{D) 1.33:} Confunde con el índice del agua; no corresponde a los datos del enunciado.
\end{itemize}
\end{respuestabox}

\begin{respuestabox}
\subsection*{Pregunta 5}
\textcolor{verdecorrecto}{\textbf{Respuesta correcta: B) 0.376}}

\textbf{Fórmula utilizada:} $n_1\sin\theta_1 = n_2\sin\theta_2$

\textbf{Sustitución:}
$$\sin\theta_2 = \frac{n_1}{n_2}\sin\theta_1 = \frac{1}{1.33} \times \sin 30^\circ = \frac{1}{1.33} \times 0.5 = 0.376$$

Como $\sin\theta_2 < \sin\theta_1 = 0.5$, el rayo se acerca a la normal al entrar al agua ($n_2 > n_1$). \checkmark

\textbf{Respuestas incorrectas:}
\begin{itemize}
    \item \textbf{A) 0.665:} Error de calcular $n_2 \times \sin\theta_1 = 1.33 \times 0.5$; invierte el cociente.
    \item \textbf{C) 1.33:} Reporta el índice de refracción del agua en lugar de calcular $\sin\theta_2$.
    \item \textbf{D) 0.500:} Reporta $\sin 30^\circ$ sin aplicar la ley de Snell.
\end{itemize}
\end{respuestabox}

\begin{respuestabox}
\subsection*{Pregunta 6}
\textcolor{verdecorrecto}{\textbf{Respuesta correcta: C) 30 cm}}

\textbf{Fórmula utilizada:} $\dfrac{1}{f} = \dfrac{1}{d_o} + \dfrac{1}{d_i}$

\textbf{Sustitución:}
$$\frac{1}{d_i} = \frac{1}{f} - \frac{1}{d_o} = \frac{1}{20} - \frac{1}{60} = \frac{3-1}{60} = \frac{2}{60}$$
$$d_i = 30 \text{ cm}$$

Como $d_i > 0$, la imagen es real y se forma frente al espejo.

\textbf{Respuestas incorrectas:}
\begin{itemize}
    \item \textbf{A) 15 cm:} Error de calcular $\frac{1}{d_i} = \frac{1}{20} - \frac{1}{20}$; usa $d_o = 20$ en lugar de 60.
    \item \textbf{B) 80 cm:} Error de sumar $\frac{1}{d_i} = \frac{1}{20} + \frac{1}{60}$ en lugar de restar.
    \item \textbf{D) 40 cm:} Error de calcular $d_i = d_o - f = 60 - 20$; confunde la ecuación con una resta directa.
\end{itemize}
\end{respuestabox}
\newpage
\begin{respuestabox}
\subsection*{Pregunta 7}
\textcolor{verdecorrecto}{\textbf{Respuesta correcta: C) $m = -0.5$; imagen invertida y reducida.}}

\textbf{Fórmula utilizada:} $m = -\dfrac{d_i}{d_o}$

\textbf{Sustitución:}
$$m = -\frac{30}{60} = -0.5$$

$m < 0$: imagen \textbf{invertida}. $|m| = 0.5 < 1$: imagen \textbf{reducida} (mide la mitad del objeto).

\textbf{Respuestas incorrectas:}
\begin{itemize}
    \item \textbf{A) $m = +0.5$:} Olvida el signo negativo en la fórmula del aumento.
    \item \textbf{B) $m = -2$:} Invierte el cociente: calcula $-d_o/d_i = -60/30$ en lugar de $-d_i/d_o$.
    \item \textbf{D) $m = +2$:} Invierte el cociente y además omite el signo negativo.
\end{itemize}
\end{respuestabox}

\begin{respuestabox}
\subsection*{Pregunta 8}
\textcolor{verdecorrecto}{\textbf{Respuesta correcta: C) 18 cm}}

\textbf{Fórmula utilizada:} $\dfrac{1}{f} = \dfrac{1}{d_o} + \dfrac{1}{d_i}$

\textbf{Sustitución:}
$$\frac{1}{d_i} = \frac{1}{12} - \frac{1}{36} = \frac{3-1}{36} = \frac{2}{36}$$
$$d_i = 18 \text{ cm}$$

Imagen real al otro lado de la lente, $m = -18/36 = -0.5$: invertida y reducida.

\textbf{Respuestas incorrectas:}
\begin{itemize}
    \item \textbf{A) 24 cm:} Error de calcular $d_o - f = 36 - 12 = 24$; confunde la ecuación con una resta.
    \item \textbf{B) 9 cm:} Error de calcular $\frac{f \cdot d_o}{f + d_o} = \frac{12 \times 36}{48} = 9$; usa suma en denominador en lugar de resta.
    \item \textbf{D) 48 cm:} Error de calcular $f + d_o = 12 + 36$; confunde la operación.
\end{itemize}
\end{respuestabox}

\begin{respuestabox}
\subsection*{Pregunta 9}
\textcolor{verdecorrecto}{\textbf{Respuesta correcta: C) A 30 cm al mismo lado del objeto; virtual y amplificada.}}

\textbf{Fórmula utilizada:} $\dfrac{1}{f} = \dfrac{1}{d_o} + \dfrac{1}{d_i}$

\textbf{Sustitución:}
$$\frac{1}{d_i} = \frac{1}{15} - \frac{1}{10} = \frac{2-3}{30} = -\frac{1}{30}$$
$$d_i = -30 \text{ cm}$$

$d_i < 0$: imagen \textbf{virtual}, del mismo lado que el objeto. $m = -(-30)/10 = +3$: imagen \textbf{derecha y amplificada} 3 veces. Este es el principio de la lupa.

\textbf{Respuestas incorrectas:}
\begin{itemize}
    \item \textbf{A):} $d_i = +30$ implicaría imagen real, pero el signo negativo indica imagen virtual al mismo lado.
    \item \textbf{B):} Confunde la posición de la imagen con la del objeto, y describe una imagen reducida siendo que $|m|=3>1$.
    \item \textbf{D):} Siempre se forma imagen; cuando el objeto está entre $F$ y la lente, la imagen es virtual y amplificada.
\end{itemize}
\end{respuestabox}
\newpage
\begin{respuestabox}
\subsection*{Pregunta 10}
\textcolor{verdecorrecto}{\textbf{Respuesta correcta: A) $-15$ cm (virtual, mismo lado que el objeto)}}

\textbf{Fórmula utilizada:} $\dfrac{1}{f} = \dfrac{1}{d_o} + \dfrac{1}{d_i}$

\textbf{Sustitución:}
$$\frac{1}{d_i} = \frac{1}{f} - \frac{1}{d_o} = \frac{1}{-20} - \frac{1}{60} = \frac{-3-1}{60} = -\frac{4}{60}$$
$$d_i = -15 \text{ cm}$$

$d_i < 0$: imagen \textbf{virtual}. La lente divergente siempre produce imágenes virtuales, derechas y reducidas. $m = -(-15)/60 = +0.25$.

\textbf{Respuestas incorrectas:}
\begin{itemize}
    \item \textbf{B) $+15$ cm:} Olvida que $f$ es negativa para lente divergente; calcula $1/20 - 1/60 = 2/60$, $d_i = +30$... o bien toma $|d_i|$ sin el signo.
    \item \textbf{C) $-30$ cm:} Error de calcular $\frac{1}{d_i} = \frac{1}{-20} - \frac{1}{60}$ incorrectamente como $-1/60 - 1/60 = -2/60$.
    \item \textbf{D) $+60$ cm:} Confunde la posición del objeto con la de la imagen; una lente divergente nunca produce imágenes reales.
\end{itemize}
\end{respuestabox}
